\chapter[\emj{🔍}~Binary Search (Búsqueda Binaria)]{\emj{🔍}~Binary Search (Búsqueda Binaria)}

\begin{dsaquees}

El juego de adivinar el número 🔢 entre 1 y 100. 

Dices "50" y te dicen: "¡Más alto!". 
Ahora sabes que el número NO está entre 1 y 50. ¡Acabas de descartar 
la MITAD de las posibilidades en un solo paso! 

Luego dices "75", te dicen: "¡Más bajo!". 
Ahora sabes que está entre 51 y 74. Vuelves a descartar la mitad.

En solo 7 intentos puedes adivinar cualquier número entre 128 opciones. 
Es increíblemente rápido porque el problema se encoge 
exponencialmente en cada paso.

\end{dsaquees}

\begin{dsaotraforma}

Binary Search es un algoritmo de búsqueda eficiente que funciona sobre
un conjunto de datos ORDENADO. 

En lugar de revisar uno por uno, siempre revisa el elemento del medio. 
Si el objetivo es menor al medio, busca en la mitad izquierda. Si es 
mayor, busca en la derecha. 

Este enfoque de "Divide y Vencerás" reduce la complejidad de tiempo
a O(log n), lo cual es ideal para datasets masivos.

\end{dsaotraforma}

\begin{dsadiagrama}
\begin{verbatim}
Buscar 23 en array ordenado:
[1, 3, 5, 7, 9, 11, 13, 15, 17, 19, 21, 23, 25]
 idx: 0  1  2  3  4   5   6   7   8   9  10  11  12

Paso 1: izq=0, der=12, mid=6.
        arr[6] es 13. ¿13 < 23? SÍ.
        Nuevo rango: [7...12] (la derecha)

Paso 2: izq=7, der=12, mid=9.
        arr[9] es 19. ¿19 < 23? SÍ.
        Nuevo rango: [10...12]

Paso 3: izq=10, der=12, mid=11.
        arr[11] es 23. ¿23 == 23? ¡SÍ! 🎉 🎉

Solo 3 pasos para 13 elementos.
\end{verbatim}
\end{dsadiagrama}

\begin{lstlisting}
Requisito: El array DEBE estar ordenado de menor a mayor.

Pasos:
1. Define dos punteros: `izq` al inicio y `der` al final.
2. Mientras `izq <= der`:
   - Calcula el punto medio: `mid = izq + (der - izq) / 2`.
   - Si `arr[mid] == objetivo`, ¡terminaste! Devuelve el índice.
   - Si `arr[mid] < objetivo`, mueve `izq = mid + 1`.
   - Si `arr[mid] > objetivo`, mueve `der = mid - 1`.
3. Si sales del loop sin encontrarlo, devuelve -1.

📝 PSEUDOCÓDIGO
──────────────────────────────────────────────────────────────

función busquedaBinaria(arr, x):
    izq = 0
    der = arr.tamaño - 1
    MIENTRAS izq <= der:
        mid = izq + (der - izq) / 2
        SI arr[mid] == x: return mid
        SI arr[mid] < x:  izq = mid + 1
        SINO:             der = mid - 1
    return -1

💻 CÓDIGO C++
──────────────────────────────────────────────────────────────

#include <iostream>
#include <vector>
using namespace std;

int binarySearch(const vector<int>& arr, int target) {
    int left = 0;
    int right = arr.size() - 1;

    while (left <= right) {
        // Usamos esta fórmula para evitar overflow de enteros
        int mid = left + (right - left) / 2;

        if (arr[mid] == target) return mid;
        
        if (arr[mid] < target)
            left = mid + 1; // Buscar en la derecha
        else
            right = mid - 1; // Buscar en la izquierda
    }
    return -1;
}

int main() {
    vector<int> sorted_data = {1, 3, 5, 7, 9, 11, 13, 15, 17, 19};
    int target = 15;

    int res = binarySearch(sorted_data, target);
    if(res != -1) cout << "Encontrado en idx: " << res << endl;
    
    return 0;
}

⏱️ COMPLEJIDAD (Big-O)
──────────────────────────────────────────────────────────────

  Caso        │ Tiempo        │ Razón
  ────────────┼───────────────┼────────────────────────────────
  Best        │ O(1)          │ El elemento está justo al medio
  Average     │ O(log n)      │ Se divide el problema a la mitad
  Worst       │ O(log n)      │ No existe o queda 1 solo elemento
  Espacio     │ O(1)          │ Versión iterativa
\end{lstlisting}

\begin{dsaentrevistas}

✅ Off-by-one errors: El bug más común es usar `left < right` en lugar
   de `left <= right`. ¡Ten cuidado!

💡 Overflow Bug: NUNCA uses `mid = (left + right) / 2`. Si left y right 
   son muy grandes, su suma puede superar el límite de un `int`. 
   Usa siempre: `left + (right - left) / 2`.

⚠️ Rotated Sorted Array: Un problema clásico de Google es buscar en un
   array ordenado que fue "rotado" (ej: [4,5,6,1,2,3]). Sigue siendo
   Binary Search pero con lógica extra.

🔑 Binary Search on Answer: No solo sirve para buscar en arrays. Si te
   preguntan "¿cuál es el mínimo valor de X que cumple esta condición?",
   puedes hacer Binary Search sobre el RANGO de posibles valores de X.

\end{dsaentrevistas}

\begin{dsaresumen}

• Solo funciona con datos ordenados.
• Divide el espacio de búsqueda a la mitad en cada paso.
• Tiempo logarítmico O(log n): Súper veloz.
• Versión iterativa usa O(1) memoria.
• Previene overflow usando la fórmula correcta de `mid`.
• Es una de las herramientas más poderosas para optimizar búsquedas.
════════════════════════════════════════════════════════════════

\end{dsaresumen}


